\documentclass[10pt,a4paper]{amsart}
\usepackage[margin=19mm]{geometry}
\usepackage{amsmath,amssymb,mathtools}
\usepackage[scr=rsfs,cal=euler]{mathalfa}
\usepackage{xfrac}
\usepackage[unicode,hidelinks,bookmarks=false]{hyperref}
\newcommand{\ie}{i.e.\ }
\newcommand{\cf}{cf.\ }
\newcommand{\resp}{respectively\ }
\newcommand{\Subsection}{\subsection}
\providecommand{\nobreakdash}{\nobreak\hbox{-}}
\setlength{\emergencystretch}{2em}
\multlinegap0pt
\usepackage{tikz}
\usetikzlibrary{cd}
\usepackage{bm}
\usepackage{amssymb}
\newtheorem{proposition}{Proposition}
\title{Free Reductive Lie Algebra Pairs of Lie-Yamaguti algebras}
\author{Sa\"id Benayadii, Martin Bordemann, Friedrich Wagemann}
\date{Source: arXiv:2606.06422v1 (CC BY 4.0)}
\begin{document}
\maketitle

\noindent\emph{Source and licence.} Free Reductive Lie Algebra Pairs of Lie-Yamaguti algebras. Version arXiv:2606.06422v1 (CC BY 4.0), distributed by arXiv under the Creative Commons Attribution 4.0 International licence, http://creativecommons.org/licenses/by/4.0/, as recorded in the arXiv licence metadata for this exact version and shown on the abstract page https://arxiv.org/abs/2606.06422v1. Full licence notice: \texttt{licenses/arxiv-2606.06422-CC-BY-4.0.txt}.\par\smallskip

\noindent\textit{Editorial scope.} Counted unit: the article's proposition proposition (source0.tex lines 464--482). The article's complete original proof of it is delivered with it (source0.tex lines 483--502, one \texttt{proof} environment of 20 lines). No mathematics is written, repaired or re-derived: the statement and the proof are the article's own lines, in the article's own notation, and no retained line is substituted or re-flowed.  Retained uncounted as the source-local context the proof needs: lines 409--416.  Nothing was omitted from any retained span, so the package carries no omission record.  The retained lines cite 1 works, whose entries are transcribed verbatim from the article's own bibliography source in the e-print and delivered with short editorial labels recorded in the item's reference mapping.  The retained lines contain the article's own diagram code, which the wrapper typesets with the standard tikz packages; no image file is delivered and the package declares no asset dependency.  Editorial formatting only, in addition: an AMS \texttt{amsart} preamble that restates the article's own declarations and macros used by the retained lines, namely the environments \texttt{proposition} on separate counters, together with the standard packages those lines need.  The retained environments print in this document as Proposition 1 in order of appearance, whereas the article prints the same items with the numbering of its own full document.  The complete native source of the article as distributed in the arXiv e-print for this exact version is bundled unchanged as \texttt{sources/202188/source0.tex}.
   \begin{equation}\label{EqCompMMComponentsMorph}
   	\forall~z_1,z_2\in\mathfrak{m}:~~
   \phi_\mathfrak{h}	\left([z_1,z_2]_\mathfrak{h}\right)
   =\left[\phi_\mathfrak{m}(z_1),\phi_\mathfrak{m}(z_1)\right]_{\mathfrak{h}'}
   ~~\mathrm{and}~~
   \phi_\mathfrak{m}	\left([z_1,z_2]_\mathfrak{m}\right)
   =\left[\phi_\mathfrak{m}(z_1),\phi_\mathfrak{m}(z_1)\right]_{\mathfrak{m}'} .   
   \end{equation}

\begin{proposition} With the above notation: $\mathfrak{i}(\mathfrak{g})$ is an
	ideal of $\mathfrak{g}$ and equal to the minimal transitive subalgebra of 
	$\mathfrak{g}$ w.r.t.~$\mathfrak{m}$. Moreover,
  the subclass $\mathrm{m}\mathcal{RLP}$ of all \textit{$\mathfrak{m}$-generated}
  RLA pairs is a full subcategory of $\mathcal{RLP}$, and there is the following
  adjunction of functors
  \begin{equation}
  	\begin{tikzcd}
  		 \mathrm{m}\mathcal{RLP}
  		\arrow[r, yshift=0.7ex, "\mathbf{J}"]
  		&\mathcal{RLP}\arrow[l, yshift=-0.7ex, 
  		"\boldsymbol{\mathfrak{i}}"]
  	\end{tikzcd}
  \end{equation}
  where $\mathbf{J}$ is the inclusion functor and $\boldsymbol{\mathfrak{i}}$
  is the functor assigning to $(\mathfrak{g},\mathfrak{h},\mathfrak{m})$
  the ideal $\mathfrak{i}(\mathfrak{g})$. The unit of the adjunction is an isomorphism and the counit is the injection
  $\mathfrak{i}(\mathfrak{g})\to \mathfrak{g}$ whence $m\mathcal{RLP}$ is thus a coreflective subcategory of $\mathcal{RLP}$, see \cite[p.~91]{Mac98}.
\end{proposition}

\begin{proof}
It is a routine check upon using $\mathfrak{h}$- and $\mathfrak{m}$-components
that $\mathfrak{i}(\mathfrak{g})$ is an ideal of $\mathfrak{g}$ which obviously is
transitive. Since every transitive subalgebra contains $\mathfrak{m}$ and 
$[\mathfrak{m},\mathfrak{m}]$, the ideal $\mathfrak{i}(\mathfrak{g})$ is contained in
every transitive subalgebra, and thus the first statement is clear. Next,
the fact that
\begin{equation}\label{EqCompHCompOfLieIdealI}
	\mathfrak{h}\cap \mathfrak{i}(\mathfrak{g})
	=K\mathrm{span}\big\{[z_1,z_2]_\mathfrak{h}
	~\big|~z_1,z_2\in\mathfrak{m}\big\}
\end{equation}
shows that $\big(\mathfrak{i}(\mathfrak{g}),\mathfrak{i}(\mathfrak{g})\cap \mathfrak{h},\mathfrak{m}\big)$ is in $\mathrm{m}\mathcal{RLP}$ and
implies that each morphism $\mathfrak{g}\to\mathfrak{g}'$ in $\mathcal{RLP}$ restricts to a morphism of RLA pairs $\mathfrak{i}(\mathfrak{g})\to
\mathfrak{i}(\mathfrak{g}')$, see (\ref{EqCompMMComponentsMorph}), whence
$\boldsymbol{\mathfrak{i}}$ is a well-defined functor. The adjunction properties are easily checked upon observing that for any
$\mathfrak{m}$-generated RLA pair $\mathfrak{g}$ and each morphism
$\mathfrak{g}\to\mathfrak{g}'$ of RLA pairs automatically corestricts to
$\mathfrak{i}(\mathfrak{g}')\subset \mathfrak{g}'$. 
\end{proof}

\begin{thebibliography}{99}
\bibitem[Mac98]{Mac98} Mac Lane, S.: \textit{Categories for the Working Mathematician}. 2nd ed., Springer, New York, 1998.
\end{thebibliography}
\end{document}
